In Table~\ref{tab:linear:algebra}, 
some finite field activities regarding linear algebra are shown.
\orbitertableofcommands{finite_field_activities_linear_algebra}
For instance, the command

\sourcecode{RREF}

computes the RREF form of the 
matrix 
$$
\left[
\begin{array}{ccccc}
1&1&1&1&0\\
1&1&0&0&1\\
\end{array}
\right]
$$
over $\bbF_2.$
The output is the matrix 
\orbiteroutput{GRAPHICS/WRAPPERS/rref_1}

\bigskip

The \verb'-RREF' command produces a latex log of the steps. 
This can be used to follow the algorithm along. 
For a somewhat longer example, 
consider the Vandermonde matrix over the field $\bbF_7$.
Suppose we want to compute the inverse matrix directly. 
We can use the following command to do so. 
Notice how we first create the matrix and an identity matrix next to it. 
After that we apply the \verb'-RREF' command:

\sourcecodevariable{V7_VANDERMONDE_EXTENDED}


\sourcecode{RREF_V7}

The following (shortened) output is produced. 
Observe how the inverse matrix 
appears in the second half once the \verb'-RREF' algorithm is finished: 
\orbiteroutput{GRAPHICS/WRAPPERS/rref_long}
The inverse matrix agrees with the output 
obtained in Section~\ref{sec:finite:fields}.



\bigskip


Another task is computing the nullspace of a matrix.
The command 

\sourcecode{nullspace}

computes the right nullspace of the matrix from the first example.
The output is the matrix
\orbiteroutput{GRAPHICS/WRAPPERS/rref_3}

\bigskip


Orbiter can compute eigenvalues 
and eigenvectors of matrices over finite fields. 
For instance, the command 

\sourcecode{eigenstuff}

computes all eigenvectors and eigenvalues of the matrix
$$
\left[
\begin{array}{cccc}
	0&1&0&2\\
	0&1&2&1\\
	4&2&3&1\\
	2&0&4&3\\
\end{array}
\right]
$$
over $\bbF_5.$

\bigskip

Orbiter can produce a list of all conjugacy classes 
of endomorphisms of $\bbF_q^d$ by means of their rational normal forms. 
For instance

\sourcecode{classes_GL_3_2}

produces a list of all conjugacy classes of $\GL(3,2).$
There are 6 of them.
The report includes the order of 
the centralizer and the order of the conjugacy class. 
The order of the centralizer is computed 
using Kung's formula~\cite{Kung81}.
This command relies on the Orbiter catalogue of irreducible polynomials.
For an introduction to the rational normal 
form of endomorphisms, see~\cite{Lueneburg1987}.
\orbiteroutput{GRAPHICS/WRAPPERS/conjugacy_classes_GL_3_2}
